% !TEX root = ../main.tex
\corrigesChapitre{4}{Avancement et bilan de matière}
\corrigesSouples

\begin{blocLivre}{cor-4-1}
\begin{corrige}{1}{Compter des molécules}
\demarche{Chaque fois que la réaction se produit, elle consomme 2 molécules de dihydrogène et 1 de dioxygène.}
\begin{questions}
\item \res{6} molécules de dihydrogène et \res{4} molécules de dioxygène.
\item Le dihydrogène permet $\dfrac{6}{2}=3$ réactions, le dioxygène 4 : la réaction se produit au plus \res{3 \text{ fois}}. Le
\res{\text{dihydrogène}} est entièrement consommé ; il reste une molécule de dioxygène.
\item État initial : 12 atomes H et 8 atomes O. État final : 6 molécules d'eau (12 H, 6 O) et une molécule de dioxygène
(2 O), soit \res{12 \text{ H et } 8 \text{ O}}.
\item Les nombres d'entités sont gigantesques : on les compte par paquets de $\num{6.02e23}$, c'est-à-dire en
\res{\text{moles}}.
\end{questions}
\schema[corrige]{04-particules}
\end{corrige}
\end{blocLivre}
\begin{blocLivre}{cor-4-2}
\begin{corrige}{2}{Le fer et le soufre}
\demarche{Les nombres stœchiométriques valent 1 : le réactif limitant est celui dont la quantité initiale est la plus petite.}
\begin{questions}
\item $m(\ce{Fe})=0{,}100\times55{,}8=\res{\qty{5,58}{g}}$ et $m(\ce{S})=0{,}080\times32{,}1=\res{\qty{2,57}{g}}$.
\item Tableau d'avancement :
\begin{avancement}{3}
\tavEquation{\ce{Fe} & $+\;\ce{S}$ & $\longrightarrow\;\ce{FeS}$}
\tavTitres{3}
initial & $0$ & $0{,}100$ & $0{,}080$ & $0$ \\
en cours & $x$ & $0{,}100-x$ & $0{,}080-x$ & $x$ \\
final & $x_\mathrm{max}=0{,}080$ & $0{,}020$ & $0$ & $0{,}080$ \\
\end{avancement}
\item $\res{x_\mathrm{max}=\qty{0,080}{mol}}$ ; le \res{\text{soufre}} est le réactif limitant.
\item $m(\ce{Fe})=0{,}020\times55{,}8=\res{\qty{1,12}{g}}$ et $m(\ce{FeS})=0{,}080\times87{,}9=\res{\qty{7,03}{g}}$. On retrouve
$1{,}12+7{,}03=\qty{8,15}{g}=5{,}58+2{,}57$ : la masse est conservée.
\end{questions}
\end{corrige}
\end{blocLivre}
\begin{blocLivre}{cor-4-3}
\begin{corrige}{3}{Réactif limitant ou mélange stœchiométrique ?}
\demarche{On compare les quotients $\dfrac{n_0}{a}$ : le plus petit donne $x_\mathrm{max}$ et désigne le réactif limitant.}
\begin{questions}
\item A : 0,10 et 0,10 ; B : 0,10 et 0,20 ; C : 0,15 et 0,10 (en \unit{mol}).
\item A : $x_\mathrm{max}=\qty{0,10}{mol}$, les deux réactifs sont épuisés ; B : $x_\mathrm{max}=\qty{0,10}{mol}$, \res{\text{magnésium limitant}} ;
C : $x_\mathrm{max}=\qty{0,10}{mol}$, \res{\text{dioxygène limitant}}.
\item Le mélange \res{\text{A}} est stœchiométrique : $\dfrac{n_0(\ce{Mg})}{2}=n_0(\ce{O2})$, soit $n_0(\ce{Mg})=2\,n_0(\ce{O2})$.
\item Mélange B : $n(\ce{Mg})=0$, $n(\ce{O2})=0{,}20-0{,}10=\res{\qty{0,10}{mol}}$ et $n(\ce{MgO})=2\times0{,}10=\res{\qty{0,20}{mol}}$.
\end{questions}
\retiens{Pour trouver le réactif limitant, on compare les quotients $\dfrac{n_0}{a}$ de chaque réactif : le plus petit donne $x_\mathrm{max}$.
Comparer directement les quantités initiales conduit à une erreur dès que les nombres stœchiométriques diffèrent.}
\end{corrige}
\end{blocLivre}
\begin{blocLivre}{cor-4-4}
\begin{corrige}{4}{Brûler du propane}
\demarche{L'équation indique qu'une mole de propane consomme 5 moles de dioxygène et produit 3 moles de dioxyde de carbone.}
\begin{questions}
\item $n=\dfrac{8{,}8}{44{,}0}=\res{\qty{0,20}{mol}}$.
\item $n(\ce{O2})=5\times0{,}20=\res{\qty{1,00}{mol}}$, soit $V=1{,}00\times24{,}0=\res{\qty{24,0}{L}}$.
\item $V_\text{air}=\dfrac{24{,}0}{0{,}21}\approx\res{\qty{114}{L}}$.
\item $n(\ce{CO2})=3\times0{,}20=\qty{0,60}{mol}$, $m=0{,}60\times44{,}0=\res{\qty{26,4}{g}}$.
\end{questions}
\end{corrige}
\end{blocLivre}
\begin{blocLivre}{cor-4-5}
\begin{corrige}{5}{Une transformation non totale}
\demarche{L'avancement final se lit sur la quantité d'un produit dans l'état final ; on le compare à $x_\mathrm{max}$.}
\begin{questions}
\item Les deux réactifs, de nombres stœchiométriques 1, sont en quantités égales : $\res{x_\mathrm{max}=\qty{1,0}{mol}}$.
\begin{avancement}{4}
\tavEquation{\ce{CH3COOH} & $+\;\ce{C2H6O}$ & $\longrightarrow\;\ce{C4H8O2}$ & $+\;\ce{H2O}$}
\tavTitres{4}
initial & $0$ & $1{,}0$ & $1{,}0$ & $0$ & $0$ \\
en cours & $x$ & $1{,}0-x$ & $1{,}0-x$ & $x$ & $x$ \\
final & $x_\mathrm{f}=0{,}67$ & $0{,}33$ & $0{,}33$ & $0{,}67$ & $0{,}67$ \\
\end{avancement}
\item $n_\mathrm{f}(\text{ester})=x_\mathrm{f}$, donc $\res{x_\mathrm{f}=\qty{0,67}{mol}}$.
\item $x_\mathrm{f}<x_\mathrm{max}$ : la transformation est \res{\text{non totale}} ; les deux réactifs restent présents.
\item \qty{0,33}{mol} d'acide, \qty{0,33}{mol} d'éthanol, \qty{0,67}{mol} d'ester et \qty{0,67}{mol} d'eau ; $m=0{,}67\times88{,}0\approx\res{\qty{59}{g}}$.
\end{questions}
\retiens{L'avancement final se mesure sur l'état final réel ; l'avancement maximal se calcule en supposant la
transformation totale. Les comparer permet de savoir si la transformation est totale.}
\end{corrige}
\end{blocLivre}
\begin{blocLivre}{cor-4-6}
\begin{corrige}{6}{Un programme pour l'état final}
\demarche{Chaque quotient $\dfrac{n_0}{a}$ est l'avancement qui épuiserait un réactif ; la transformation s'arrête au premier atteint.}
\begin{questions}
\item \code{candidats} vaut \code{[0.04, 0.05]}. La transformation s'arrête dès qu'un réactif est épuisé, donc pour la
\res{\text{plus petite}} de ces valeurs.
\item \code{candidats} devient \code{[0.04, 0.03]} : le programme affiche \code{xmax = 0.03 mol}, \code{CH4 0.01 mol},
\code{O2 0.0 mol}, \code{CO2 0.03 mol} et \code{H2O 0.06 mol}. Le \res{\text{dioxygène}} est désormais limitant.
\item \code{reactifs = \{'Fe': [1, 0.100], 'S': [1, 0.080]\}} et \code{produits = \{'FeS': [1, 0.0]\}}. Affichage :
\code{xmax = 0.08 mol}, \code{Fe 0.02 mol}, \code{S 0.0 mol} et \code{FeS 0.08 mol}, comme dans l'exercice~2.
\item Dans la première boucle : \code{if n0 / a == xmax: print(espece, 'est limitant')}. Le test compare exactement les
valeurs dont \code{xmax} est le minimum.
\end{questions}
\end{corrige}
\end{blocLivre}
\begin{blocLivre}{cor-4-7}
\begin{corrige}{7}{Le coussin gonflable de sécurité}
\demarche{Avec un nombre stœchiométrique 2 pour l'azoture, $x_\mathrm{max}=\dfrac{n_0(\ce{NaN3})}{2}$ ; chaque espèce s'en déduit.}
\begin{questions}
\item $M(\ce{NaN3})=23{,}0+3\times14{,}0=\qty{65,0}{g.mol^{-1}}$, $n=\dfrac{130}{65{,}0}=\res{\qty{2,00}{mol}}$ et $\res{x_\mathrm{max}=\qty{1,00}{mol}}$.
\item $n(\ce{N2})=3x_\mathrm{max}=\qty{3,00}{mol}$ et $V=3{,}00\times24{,}0=\res{\qty{72,0}{L}}$ : c'est \res{\text{suffisant}} pour gonfler le coussin, le surplus
s'échappant par des évents.
\item $n(\ce{Na})=2x_\mathrm{max}=\qty{2,00}{mol}$, soit $m=2{,}00\times23{,}0=\res{\qty{46,0}{g}}$.
\item $n(\ce{KNO3})=\dfrac{2}{10}\times2{,}00=\qty{0,400}{mol}$ ; $M(\ce{KNO3})=\qty{101,1}{g.mol^{-1}}$, d'où $m=\res{\qty{40,4}{g}}$.
\end{questions}
\piege{Prendre $x_\mathrm{max}=n_0(\ce{NaN3})$ : le nombre stœchiométrique de l'azoture vaut 2, donc $x_\mathrm{max}=\dfrac{n_0}{2}$. Les quantités de
produits s'en déduisent : $2x_\mathrm{max}$ de sodium et $3x_\mathrm{max}$ de diazote.}
\end{corrige}
\end{blocLivre}
\begin{blocLivre}{cor-4-8}
\begin{corrige}{8}{Souder des rails}
\demarche{Mélange stœchiométrique : $\dfrac{n_0(\ce{Al})}{2}=\dfrac{n_0(\ce{Fe2O3})}{1}$ ; les deux réactifs sont épuisés en même temps.}
\begin{questions}
\item $n(\ce{Al})=\dfrac{2700}{27{,}0}=\res{\qty{100}{mol}}$ ; $M(\ce{Fe2O3})=\qty{159,6}{g.mol^{-1}}$ et $n(\ce{Fe2O3})=\dfrac{7980}{159{,}6}=\res{\qty{50,0}{mol}}$.
\item $\dfrac{100}{2}=\qty{50,0}{mol}=\dfrac{50{,}0}{1}$ : le mélange est \res{\text{stœchiométrique}} et $x_\mathrm{max}=\qty{50,0}{mol}$.
\item $n(\ce{Fe})=2x_\mathrm{max}=\qty{100}{mol}$, $m=100\times55{,}8=\res{\qty{5,58}{kg}}$.
\item $M(\ce{Al2O3})=\qty{102,0}{g.mol^{-1}}$ et $m=50{,}0\times102{,}0=\res{\qty{5,10}{kg}}$. On a bien $5{,}58+5{,}10=\qty{10,68}{kg}=2{,}70+7{,}98$.
\end{questions}
\end{corrige}
\end{blocLivre}
\begin{blocLivre}{cor-4-9}
\begin{corrige}{9}{Vinaigre et bicarbonate}
\demarche{La perte de masse est la masse de dioxyde de carbone qui s'est échappé ; une mole d'avancement libère une mole de gaz.}
\begin{questions}
\item $n(\ce{HCO3-})=\dfrac{2{,}10}{84{,}0}=\res{\qty{0,0250}{mol}}$ ; $c=\dfrac{50}{60{,}0}=\qty{0,833}{mol.L^{-1}}$ et $n(\ce{CH3COOH})=0{,}833\times0{,}0500=\res{\qty{4,17e-2}{mol}}$.
\item $\res{x_\mathrm{max}=\qty{0,0250}{mol}}$, l'\res{\text{hydrogénocarbonate}} est limitant ; on attend $0{,}0250\times44{,}0=\res{\qty{1,10}{g}}$ de dioxyde de carbone.
\item La perte de masse vaut $152{,}48-151{,}41=\qty{1,07}{g}$, d'où $x_\mathrm{f}=\dfrac{1{,}07}{44{,}0}=\res{\qty{0,0243}{mol}}$.
\item $\dfrac{x_\mathrm{f}}{x_\mathrm{max}}\approx0{,}97$ : on peut considérer la transformation comme \res{\text{totale}}. Un peu de dioxyde de carbone
reste dissous dans la solution, si bien que la perte de masse sous-estime légèrement l'avancement.
\end{questions}
\end{corrige}
\end{blocLivre}
\begin{blocLivre}{cor-4-10}
\begin{corrige}{10}{Un précipité blanc}
\demarche{Seuls les ions $\ce{Ag+}$ et $\ce{Cl-}$ réagissent ; les ions $\ce{Na+}$ et $\ce{NO3-}$, spectateurs, gardent leur quantité.}
\begin{questions}
\item $n(\ce{Ag+})=0{,}0200\times0{,}050=\res{\qty{1,00}{mmol}}$ et $n(\ce{Cl-})=0{,}0100\times0{,}080=\res{\qty{0,80}{mmol}}$.
\item Les ions chlorure sont \res{\text{limitants}} (quantités en \unit{mol}) :
\begin{avancement}{3}
\tavEquation{\ce{Ag+} & $+\;\ce{Cl-}$ & $\longrightarrow\;\ce{AgCl}$}
\tavTitres{3}
initial & $0$ & $\num{1.00e-3}$ & $\num{0.80e-3}$ & $0$ \\
en cours & $x$ & $\num{1.00e-3}-x$ & $\num{0.80e-3}-x$ & $x$ \\
final & $x_\mathrm{max}=\num{0.80e-3}$ & $\num{0.20e-3}$ & $0$ & $\num{0.80e-3}$ \\
\end{avancement}
\item $m=\num{0.80e-3}\times143{,}4=\res{\qty{0,115}{g}}$.
\item $[\ce{Ag+}]=\dfrac{\num{0.20e-3}}{0{,}0300}=\res{\qty{6,7e-3}{mol.L^{-1}}}$ ; $[\ce{NO3-}]=\dfrac{\num{1.00e-3}}{0{,}0300}=\res{\qty{3,33e-2}{mol.L^{-1}}}$ ;
$[\ce{Na+}]=\dfrac{\num{0.80e-3}}{0{,}0300}=\res{\qty{2,67e-2}{mol.L^{-1}}}$ ; les ions chlorure ont disparu.
\end{questions}
\end{corrige}
\end{blocLivre}
\begin{blocLivre}{cor-4-11}
\begin{corrige}{11}{Retrouver une quantité initiale}
\demarche{La mesure de l'état final donne une équation du premier degré en $x_\mathrm{f}$ ; le zinc épuisé donne ensuite $n_0=x_\mathrm{f}$.}
\begin{questions}
\item Avec $n_0(\ce{Cu^{2+}})=0{,}0500\times0{,}100=\qty{5,00e-3}{mol}$ :
\begin{avancement}{4}
\tavEquation{\ce{Zn} & $+\;\ce{Cu^{2+}}$ & $\longrightarrow\;\ce{Zn^{2+}}$ & $+\;\ce{Cu}$}
\tavTitres{4}
initial & $0$ & $n_0$ & $\num{5.00e-3}$ & $0$ & $0$ \\
final & $x_\mathrm{f}$ & $n_0-x_\mathrm{f}=0$ & $\num{5.00e-3}-x_\mathrm{f}$ & $x_\mathrm{f}$ & $x_\mathrm{f}$ \\
\end{avancement}
\item $n_\mathrm{f}(\ce{Cu^{2+}})=0{,}040\times0{,}0500=\qty{2,00e-3}{mol}$, donc $\num{5.00e-3}-x_\mathrm{f}=\num{2.00e-3}$ et $\res{x_\mathrm{f}=\qty{3,00e-3}{mol}}$.
\item Le zinc est épuisé : $n_0=x_\mathrm{f}=\qty{3,00e-3}{mol}$ et $m=\num{3.00e-3}\times65{,}4=\res{\qty{0,196}{g}}$.
\item $[\ce{Zn^{2+}}]=\dfrac{\num{3.00e-3}}{0{,}0500}=\res{\qty{0,060}{mol.L^{-1}}}$.
\end{questions}
\retiens{Quand l'état final est mesuré, la relation $n_\mathrm{f}=n_0-a\,x_\mathrm{f}$ devient une équation du premier degré dont l'inconnue est
l'avancement final, ou une quantité initiale.}
\end{corrige}
\end{blocLivre}
\begin{blocLivre}{cor-4-12}
\begin{corrige}{12}{Un réchaud de camping}
\demarche{Le butane a pour nombre stœchiométrique 2 : $x_\mathrm{max}=\dfrac{n_0}{2}$, puis $n(\ce{O2})=13x_\mathrm{max}$, $n(\ce{CO2})=8x_\mathrm{max}$, $n(\ce{H2O})=10x_\mathrm{max}$.}
\begin{questions}
\item $M(\ce{C4H10})=\qty{58,0}{g.mol^{-1}}$, $n=\dfrac{190}{58{,}0}=\res{\qty{3,28}{mol}}$ et $\res{x_\mathrm{max}=\qty{1,64}{mol}}$.
\item $n(\ce{O2})=13\times1{,}638=\qty{21,3}{mol}$, $V=\res{\qty{511}{L}}$ ; $V_\text{air}=\dfrac{511}{0{,}21}\approx\res{\qty{2,43}{m^3}}$.
\item $m(\ce{CO2})=8\times1{,}638\times44{,}0\approx\res{\qty{577}{g}}$ et $m(\ce{H2O})=10\times1{,}638\times18{,}0\approx\res{\qty{295}{g}}$.
\item Brûler toute la cartouche consommerait \qty{0,51}{m^3} de dioxygène, soit environ \qty{85}{\percent} de celui de la tente. Bien avant,
la teneur en dioxygène baisse : le \res{\text{dioxygène devient le réactif limitant}}, la combustion devient incomplète et
produit du monoxyde de carbone, tandis que les occupants manquent d'air.
\end{questions}
\end{corrige}
\end{blocLivre}
\begin{blocLivre}{cor-4-13}
\begin{corrige}{13}{Fabriquer de l'ammoniac}
\demarche{On lit $x_\mathrm{f}$ sur la quantité d'ammoniac formé : $n(\ce{NH3})=2x_\mathrm{f}$.}
\begin{questions}
\item $\dfrac{1{,}00}{1}=\dfrac{3{,}00}{3}$ : le mélange est \res{\text{stœchiométrique}} et $\res{x_\mathrm{max}=\qty{1,00}{mol}}$.
\item $2x_\mathrm{f}=0{,}30$, donc $\res{x_\mathrm{f}=\qty{0,15}{mol}}$ :
\begin{avancement}{3}
\tavEquation{\ce{N2} & $+\;3\,\ce{H2}$ & $\longrightarrow\;2\,\ce{NH3}$}
\tavTitres{3}
initial & $0$ & $1{,}00$ & $3{,}00$ & $0$ \\
en cours & $x$ & $1{,}00-x$ & $3{,}00-3x$ & $2x$ \\
final & $x_\mathrm{f}=0{,}15$ & $0{,}85$ & $2{,}55$ & $0{,}30$ \\
\end{avancement}
\item $x_\mathrm{f}<x_\mathrm{max}$ : \res{\text{non totale}}. Le mélange sortant contient \qty{0,85}{mol} de diazote, \qty{2,55}{mol} de dihydrogène et
\qty{0,30}{mol} d'ammoniac.
\item Après $k$ passages, il reste la fraction $0{,}85^k$ du diazote. Or $0{,}85^{14}\approx0{,}103$ et $0{,}85^{15}\approx0{,}087$ : il faut
\res{15 \text{ passages}}.
\end{questions}
\end{corrige}
\end{blocLivre}
\begin{blocLivre}{cor-4-14}
\begin{corrige}{14}{La lampe à carbure}
\demarche{L'acétylène formé par la première réaction est le réactif de la seconde : on enchaîne les deux bilans.}
\begin{questions}
\item $M(\ce{CaC2})=40{,}1+2\times12{,}0=\qty{64,1}{g.mol^{-1}}$ et $n=\dfrac{100}{64{,}1}=\res{\qty{1,56}{mol}}$.
\item $n(\ce{H2O})=2n=\qty{3,12}{mol}$, soit $\res{\qty{56,2}{g}}$ d'eau ; $n(\ce{C2H2})=n=\qty{1,56}{mol}$, soit $V=\res{\qty{37,4}{L}}$.
\item $n(\ce{O2})=\dfrac{5}{2}\times1{,}56=\qty{3,90}{mol}$, $V=\res{\qty{93,6}{L}}$.
\item Autonomie : $\dfrac{37{,}4}{12}\approx\res{\qty{3,1}{h}}$, soit environ trois heures d'éclairage par charge.
\end{questions}
\end{corrige}
\end{blocLivre}
\begin{blocLivre}{cor-4-15}
\begin{corrige}{15}{La paille de fer qui brûle}
\demarche{La pente de chaque droite vaut $-a$ pour un réactif et $+c$ pour un produit : elle donne les nombres stœchiométriques.}
\begin{questions}
\item La droite \res{\text{c}} part de zéro et monte : c'est l'oxyde magnétique, seul produit.
\item Pentes : a, $\dfrac{0-0{,}30}{0{,}10}=-3$ ; b, $\dfrac{0{,}05-0{,}25}{0{,}10}=-2$ ; c, $+1$. Une entité $\ce{Fe3O4}$ contient 3 atomes de fer et 4 atomes
d'oxygène, soit 2 molécules de dioxygène : a est le \res{\text{fer}}, b le \res{\text{dioxygène}}, et $\res{\ce{3Fe + 2O2 -> Fe3O4}}$.
\item $n_0(\ce{Fe})=\qty{0,30}{mol}$, soit $\res{\qty{16,7}{g}}$ ; $n_0(\ce{O2})=\qty{0,25}{mol}$, soit $\res{\qty{6,0}{L}}$.
\item $\dfrac{0{,}30}{3}=0{,}10$ et $\dfrac{0{,}25}{2}=0{,}125$ : $\res{x_\mathrm{max}=\qty{0,10}{mol}}$, le \res{\text{fer}} est limitant. $M(\ce{Fe3O4})=\qty{231,4}{g.mol^{-1}}$ et
$m=\res{\qty{23,1}{g}}$.
\end{questions}
\schema[corrige]{04-ex15}
\end{corrige}
\end{blocLivre}
